\documentclass[10pt,a4paper]{article}
\PassOptionsToPackage{hyphens}{url}
\usepackage{amsmath,amssymb,amsthm,graphicx,verbatim,fontspec,hyperref}
\setmainfont{lmroman10-regular.otf}
\emergencystretch=3em
\usepackage[a4paper,margin=20mm]{geometry}
% Retained preamble for the Stacks Project chapter snapshot in this directory.
% The chapter's own preamble.tex is a copy of the shared Stacks Project preamble; the block
% below is sliced MECHANICALLY, line for line, from it. Nothing was added, removed or
% redefined; the wrapper class, the xr-hyper/hyperref wiring and the externaldocument list
% are omitted because the assembler supplies the document class.
%
% lines taken from preamble.tex: 3, 5, 9, 11, 12, 156, 157, 158, ... (51 lines total)
%
% package coverage restored from the chapter's own preamble (2026-10-13): xy, xy-2cell

\usepackage[all]{xy}
\xyoption{2cell}
\UseAllTwocells

\newenvironment{reference}{}{}
\newenvironment{slogan}{\comment}{\endcomment}
\newenvironment{history}{\comment}{\endcomment}
\theoremstyle{plain}
\newtheorem{theorem}[subsection]{Theorem}
\newtheorem{proposition}[subsection]{Proposition}
\newtheorem{lemma}[subsection]{Lemma}
\theoremstyle{definition}
\newtheorem{definition}[subsection]{Definition}
\newtheorem{example}[subsection]{Example}
\newtheorem{exercise}[subsection]{Exercise}
\newtheorem{situation}[subsection]{Situation}
\theoremstyle{remark}
\newtheorem{remark}[subsection]{Remark}
\newtheorem{remarks}[subsection]{Remarks}
\numberwithin{equation}{subsection}
\def\lim{\mathop{\mathrm{lim}}\nolimits}
\def\colim{\mathop{\mathrm{colim}}\nolimits}
\def\Spec{\mathop{\mathrm{Spec}}}
\def\Hom{\mathop{\mathrm{Hom}}\nolimits}
\def\Ext{\mathop{\mathrm{Ext}}\nolimits}
\def\SheafHom{\mathop{\mathcal{H}\!\mathit{om}}\nolimits}
\def\SheafExt{\mathop{\mathcal{E}\!\mathit{xt}}\nolimits}
\def\Sch{\mathit{Sch}}
\def\Mor{\mathop{\mathrm{Mor}}\nolimits}
\def\Ob{\mathop{\mathrm{Ob}}\nolimits}
\def\Sh{\mathop{\mathit{Sh}}\nolimits}
\def\NL{\mathop{N\!L}\nolimits}
\def\CH{\mathop{\mathrm{CH}}\nolimits}
\def\proetale{{pro\text{-}\acute{e}tale}}
\def\etale{{\acute{e}tale}}
\def\QCoh{\mathit{QCoh}}
\def\Ker{\mathop{\mathrm{Ker}}}
\def\Im{\mathop{\mathrm{Im}}}
\def\Coker{\mathop{\mathrm{Coker}}}
\def\Coim{\mathop{\mathrm{Coim}}}
\DeclareMathSymbol{\boxtimes}{\mathbin}{AMSa}{"02}
\def\QCohstack{\mathcal{QC}\!\mathit{oh}}
\def\Cohstack{\mathcal{C}\!\mathit{oh}}
\def\Spacesstack{\mathcal{S}\!\mathit{paces}}
\def\Quotfunctor{\mathrm{Quot}}
\def\Hilbfunctor{\mathrm{Hilb}}
\def\Curvesstack{\mathcal{C}\!\mathit{urves}}
\def\Polarizedstack{\mathcal{P}\!\mathit{olarized}}
\def\Complexesstack{\mathcal{C}\!\mathit{omplexes}}
\def\Pic{\mathop{\mathrm{Pic}}\nolimits}
\def\Picardstack{\mathcal{P}\!\mathit{ic}}
\def\Picardfunctor{\mathrm{Pic}}
\def\Deformationcategory{\mathcal{D}\!\mathit{ef}}
\title{Construction and universal exactness of the Gabriel–Serre quotient}
\author{The Stacks Project Authors}
\date{stacks-project:homology.tex@a04446e57ec1fbc252a871afcec7752fb2807b14; GNU FDL 1.2 or later}
\begin{document}
\maketitle
\noindent\textit{Editorial note.} Original mathematical text and complete proof retained. Native reference metadata is visible for attribution, and the original shared bibliography is supplied. External references display exact original labels without invented numbers. No proof or hypothesis repaired.\par\medskip
\noindent\textit{Allowed standard prerequisites (editorial).}\begin{itemize}
\item Abelian categories and basic homological algebra. Original localization and Serre-subcategory definitions and supporting proofs are included.. The quotient category is constructed, not supplied as a free premise.
\end{itemize}\medskip
\section{Localization}
\label{section-localization}

\noindent
In this section we note how Gabriel-Zisman localization interacts with
the additive structure on a category.

\begin{lemma}
\label{lemma-localization-preadditive}
Let $\mathcal{C}$ be a preadditive category.
Let $S$ be a left or right multiplicative system.
There exists a unique preadditive structure on
$S^{-1}\mathcal{C}$ such that the localization functor
$Q : \mathcal{C} \to S^{-1}\mathcal{C}$ is additive.
\end{lemma}

\begin{proof}
We will prove this in the case $S$ is a left multiplicative system.
The case where $S$ is a right multiplicative system is dual.
Suppose that $X, Y$ are objects of $\mathcal{C}$ and that
$\alpha, \beta : X \to Y$ are morphisms in $S^{-1}\mathcal{C}$. According to
Categories, Lemma \textnormal{[original label: \nolinkurl{categories-lemma-morphisms-left-localization}]}
we may represent these by pairs $s^{-1}f, s^{-1}g$ with common denominator $s$.
Suppose that we have a pre-additive structure on $S^{-1}\mathcal{C}$ and
that $Q$ is additive. Then
\begin{align*}
s^{-1}f+s^{-1}g
&=Q(s)^{-1}\circ Q(f)+Q(s)^{-1}\circ Q(g)\\
&=Q(s)^{-1}\circ(Q(f)+Q(g))\\
&=Q(s)^{-1}\circ Q(f+g)\\
&=s^{-1}(f+g).
\end{align*}
Hence, the pre-additive structure on $S^{-1}\mathcal{C}$ is unique if it
exists. Correspondingly, we define $\alpha + \beta$ to be the
equivalence class of $s^{-1}(f + g)$ and in the rest of the proof we show
that this is well defined and that composition is bilinear.
Once this is done it is clear that $Q$ is an additive functor.

\medskip\noindent
Let us show construction above is well defined.
An abstract way of saying this is that filtered colimits of
abelian groups agree with filtered colimits of sets and to use
Categories,
Equation (\textnormal{[original label: \nolinkurl{categories-equation-left-localization-morphisms-colimit}]}).
We can work this out in a bit more detail as follows.
Say $s : Y \to Y_1$ and $f, g : X \to Y_1$. Suppose we have a second
representation of $\alpha, \beta$ as $(s')^{-1}f', (s')^{-1}g'$ with
$s' : Y \to Y_2$ and $f', g' : X \to Y_2$. By
Categories, Remark \textnormal{[original label: \nolinkurl{categories-remark-left-localization-morphisms-colimit}]}
we can find a morphism $s_3 : Y \to Y_3$ and morphisms
$a_1 : Y_1 \to Y_3$, $a_2 : Y_2 \to Y_3$ such that
$a_1 \circ s = s_3 = a_2 \circ s'$ and also
$a_1 \circ f = a_2 \circ f'$ and $a_1 \circ g = a_2 \circ g'$.
Hence we see that $s^{-1}(f + g)$ is equivalent to
\begin{align*}
s_3^{-1}(a_1 \circ (f + g)) & =
s_3^{-1}(a_1 \circ f + a_1 \circ g) \\
& = s_3^{-1}(a_2 \circ f' + a_2 \circ g') \\
& = s_3^{-1}(a_2 \circ (f' + g'))
\end{align*}
which is equivalent to $(s')^{-1}(f' + g')$.

\medskip\noindent
Fix $s : Y \to Y'$ and $f, g : X \to Y'$ with
$\alpha = s^{-1}f$ and $\beta = s^{-1}g$ as morphisms $X \to Y$
in $S^{-1}\mathcal{C}$.
To show that composition is bilinear first consider the case of a
morphism $\gamma : Y \to Z$ in $S^{-1}\mathcal{C}$. Say $\gamma = t^{-1}h$
for some $h : Y \to Z'$ and $t : Z \to Z'$ in $S$. Using LMS2 we
choose morphisms $a : Y' \to Z''$ and $t' : Z' \to Z''$ in $S$ such
that $a \circ s = t' \circ h$. Picture
$$
\xymatrix{
& & Z \ar[d]^t \\
& Y \ar[r]^h \ar[d]^s & Z' \ar[d]^{t'} \\
X \ar[r]^{f, g} & Y' \ar[r]^a & Z''
}
$$
Then
$\gamma \circ \alpha = (t' \circ t)^{-1}(a \circ f)$ and
$\gamma \circ \beta = (t' \circ t)^{-1}(a \circ g)$.
Hence we see that $\gamma \circ (\alpha + \beta)$ is represented
by $(t' \circ t)^{-1}(a \circ (f + g)) =
(t' \circ t)^{-1}(a \circ f + a \circ g)$ which represents
$\gamma \circ \alpha + \gamma \circ \beta$.

\medskip\noindent
Finally, assume that $\delta : W \to X$ is another morphism of
$S^{-1}\mathcal{C}$. Say $\delta = r^{-1}i$ for some
$i : W \to X'$ and $r : X \to X'$ in $S$. We claim that we can find
a morphism $s' : Y' \to Y''$ in $S$ and morphisms $a'', b'' : X' \to Y''$
such that the following diagram commutes
$$
\xymatrix{
& & & Y \ar[d]^s \\
& X \ar[rr]^{f, g, f + g} \ar[d]^r & & Y' \ar[d]^{s'} \\
W \ar[r]^i & X' \ar[rr]^{a'', b'', a'' + b''} & & Y''
}
$$
Namely, using LMS2 we can first choose
$s_1 : Y' \to Y_1$, $s_2 : Y' \to Y_2$ in $S$ and
$a : X' \to Y_1$, $b : X' \to Y_2$ such that
$a \circ r = s_1 \circ f$ and $b \circ r = s_2 \circ f$.
Then using that the category $Y'/S$ is filtered (see
Categories, Remark \textnormal{[original label: \nolinkurl{categories-remark-left-localization-morphisms-colimit}]}),
we can
find a $s' : Y' \to Y''$ and morphisms $a' : Y_1 \to Y''$, $b' : Y_2 \to Y''$
such that $s' = a' \circ s_1$ and $s' = b' \circ s_2$. Setting
$a'' = a' \circ a$ and $b'' = b' \circ b$ works.
At this point we see that the compositions
$\alpha \circ \delta$ and $\beta \circ \delta$ are represented by
$(s' \circ s)^{-1}(a'' \circ i)$ and $(s' \circ s)^{-1}(b'' \circ i)$.
Hence $\alpha \circ \delta + \beta \circ \delta$ is represented
by $(s' \circ s)^{-1}(a'' \circ i + b'' \circ i) =
(s' \circ s)^{-1}((a'' + b'') \circ i)$
which by the diagram again is a representative
of $(\alpha + \beta) \circ \delta$.
\end{proof}

\begin{lemma}
\label{lemma-localization-additive}
Let $\mathcal{C}$ be an additive category.
Let $S$ be a left or right multiplicative system.
Then $S^{-1}\mathcal{C}$ is an additive category and the localization functor
$Q : \mathcal{C} \to S^{-1}\mathcal{C}$ is additive.
\end{lemma}

\begin{proof}
By Lemma \ref{lemma-localization-preadditive}
we see that $S^{-1}\mathcal{C}$ is preadditive and that $Q$ is additive.
By Lemma \textnormal{[original label: \nolinkurl{lemma-additive-additive}]} we conclude that $S^{-1}\mathcal{C}$
has a zero object and direct sums.
\end{proof}

\noindent
The following lemma describes the ``kernel''
of the localization functor in case we invert a multiplicative system.

\begin{lemma}
\label{lemma-kernel-localization}
Let $\mathcal{C}$ be an additive category. Let $S$ be a multiplicative
system. Let $X$ be an object
of $\mathcal{C}$. The following are equivalent
\begin{enumerate}
\item $Q(X) = 0$ in $S^{-1}\mathcal{C}$,
\item there exists $Y \in \Ob(\mathcal{C})$ such that
$0 : X \to Y$ is an element of $S$, and
\item there exists $Z \in \Ob(\mathcal{C})$ such that
$0 : Z \to X$ is an element of $S$.
\end{enumerate}
\end{lemma}

\begin{proof}
If (2) holds we see that $0 = Q(0) : Q(X) \to Q(Y)$ is an isomorphism.
In the additive category $S^{-1}\mathcal{C}$ this implies that $Q(X) = 0$.
Hence (2) $\Rightarrow$ (1). Similarly, (3) $\Rightarrow$ (1).
Suppose that $Q(X) = 0$. This implies that the morphism
$f : 0 \to X$ is transformed into an isomorphism in $S^{-1}\mathcal{C}$.
Hence by
Categories, Lemma \textnormal{[original label: \nolinkurl{categories-lemma-what-gets-inverted}]}
there exists a morphism $g : Z \to 0$ such that $fg \in S$. This proves
(1) $\Rightarrow$ (3). Similarly, (1) $\Rightarrow$ (2).
\end{proof}

\begin{lemma}
\label{lemma-localization-abelian}
Let $\mathcal{A}$ be an abelian category.
\begin{enumerate}
\item If $S$ is a left multiplicative system, then
the category $S^{-1}\mathcal{A}$ has cokernels and the functor
$Q : \mathcal{A} \to S^{-1}\mathcal{A}$ commutes with them.
\item If $S$ is a right multiplicative system, then
the category $S^{-1}\mathcal{A}$ has kernels and the functor
$Q : \mathcal{A} \to S^{-1}\mathcal{A}$ commutes with them.
\item If $S$ is a multiplicative system, then the category
$S^{-1}\mathcal{A}$ is abelian and the functor
$Q : \mathcal{A} \to S^{-1}\mathcal{A}$ is exact.
\end{enumerate}
\end{lemma}

\begin{proof}
Assume $S$ is a left multiplicative system. Let $a : X \to Y$ be a morphism
of $S^{-1}\mathcal{A}$. Then $a = s^{-1}f$ for some $s : Y \to Y'$
in $S$ and $f : X \to Y'$. Since $Q(s)$ is an isomorphism we see that
the existence of $\Coker(a : X \to Y)$ is equivalent to the existence
of $\Coker(Q(f) : X \to Y')$. Since $\Coker(Q(f))$ is the
coequalizer of $0$ and $Q(f)$ we see that $\Coker(Q(f))$ is
represented by $Q(\Coker(f))$ by
Categories, Lemma \textnormal{[original label: \nolinkurl{categories-lemma-left-localization-limits}]}.
This proves (1).

\medskip\noindent
Part (2) is dual to part (1).

\medskip\noindent
If $S$ is a multiplicative system, then $S$ is both a left and a right
multiplicative system. Thus we see that $S^{-1}\mathcal{A}$ has
kernels and cokernels and $Q$ commutes with kernels and cokernels.
To finish the proof of (3) we have to show that $\Coim = \Im$ in
$S^{-1}\mathcal{A}$. Again using that any arrow in $S^{-1}\mathcal{A}$
is isomorphic to an arrow $Q(f)$ we see that the result follows
from the result for $\mathcal{A}$.
\end{proof}




\section{Jordan-H\"older}
\label{section-jordan-holder}

\noindent
The Jordan-H\"older lemma is Lemma \ref{lemma-jordan-holder}.
First we state some definitions.

\begin{definition}
\label{definition-simple}
Let $\mathcal{A}$ be an abelian category. An object $A$ of $\mathcal{A}$
is said to be {\it simple} if it is nonzero and the only subobjects
of $A$ are $0$ and $A$.
\end{definition}

\begin{definition}
\label{definition-Artinian}
Let $\mathcal{A}$ be an abelian category.
\begin{enumerate}
\item We say an object $A$ of $\mathcal{A}$ is {\it Artinian} if and only if
it satisfies the descending chain condition for subobjects.
\item We say $\mathcal{A}$ is {\it Artinian} if every object of
$\mathcal{A}$ is Artinian.
\end{enumerate}
\end{definition}

\begin{definition}
\label{definition-Noetherian}
Let $\mathcal{A}$ be an abelian category.
\begin{enumerate}
\item We say an object $A$ of $\mathcal{A}$ is {\it Noetherian} if and only if
it satisfies the ascending chain condition for subobjects.
\item We say $\mathcal{A}$ is {\it Noetherian} if every object of
$\mathcal{A}$ is Noetherian.
\end{enumerate}
\end{definition}

\begin{lemma}
\label{lemma-ses-artinian}
Let $\mathcal{A}$ be an abelian category. Let $0 \to A_1 \to A_2 \to A_3 \to 0$
be a short exact sequence of $\mathcal{A}$. Then $A_2$ is Artinian
if and only if $A_1$ and $A_3$ are Artinian.
\end{lemma}

\begin{proof}
Omitted.
\end{proof}

\begin{lemma}
\label{lemma-ses-noetherian}
Let $\mathcal{A}$ be an abelian category. Let $0 \to A_1 \to A_2 \to A_3 \to 0$
be a short exact sequence of $\mathcal{A}$. Then $A_2$ is Noetherian
if and only if $A_1$ and $A_3$ are Noetherian.
\end{lemma}

\begin{proof}
Omitted.
\end{proof}

\begin{lemma}
\label{lemma-finite-length}
Let $\mathcal{A}$ be an abelian category. Let $A$ be an object
of $\mathcal{A}$. The following are equivalent
\begin{enumerate}
\item $A$ is Artinian and Noetherian, and
\item there exists a filtration
$0 \subset A_1 \subset A_2 \subset \ldots \subset A_n = A$
by subobjects such that $A_i/A_{i - 1}$ is simple for $i = 1, \ldots, n$.
\end{enumerate}
\end{lemma}

\begin{proof}
Assume (1). If $A$ is zero, then (2) holds. If $A$ is not zero, then
there exists a smallest nonzero object $A_1 \subset A$ by the Artinian
property. Of course $A_1$ is simple. If $A_1 = A$, then we are done.
If not, then we can find $A_1 \subset A_2 \subset A$ minimal
with $A_2 \not = A_1$. Then $A_2/A_1$ is simple. Continuing in this way, we
can find a sequence $0 \subset A_1 \subset A_2 \subset \ldots $
of subobjects of $A$ such that $A_i/A_{i - 1}$ is simple. Since $A$
is Noetherian, we conclude that the process stops. Hence (2) follows.

\medskip\noindent
Assume (2). We will prove (1) by induction on $n$. If $n = 1$, then
$A$ is simple and clearly Noetherian and Artinian. If the result holds
for $n - 1$, then we use the short exact sequence
$0 \to A_{n - 1} \to A_n \to A_n/A_{n - 1} \to 0$
and Lemmas \ref{lemma-ses-artinian} and \ref{lemma-ses-noetherian}
to conclude for $n$.
\end{proof}

\begin{lemma}[Jordan-H\"older]
\label{lemma-jordan-holder}
Let $\mathcal{A}$ be an abelian category. Let $A$ be an object
of $\mathcal{A}$ satisfying the equivalent conditions of
Lemma \ref{lemma-finite-length}. Given two filtrations
$$
0 \subset A_1 \subset A_2 \subset \ldots \subset A_n = A
\quad\text{and}\quad
0 \subset B_1 \subset B_2 \subset \ldots \subset B_m = A
$$
with $S_i = A_i/A_{i - 1}$ and $T_j = B_j/B_{j - 1}$ simple objects we have
$n = m$ and there exists a permutation $\sigma$ of $\{1, \ldots, n\}$
such that $S_i \cong T_{\sigma(i)}$ for all $i \in \{1, \ldots, n\}$.
\end{lemma}

\begin{proof}
Let $j$ be the smallest index such that $A_1 \subset B_j$.
Then the map $S_1 = A_1 \to B_j/B_{j - 1} = T_j$ is an isomorphism.
Moreover, the object $A/A_1 = A_n/A_1 = B_m/A_1$
has the two filtrations
$$
0 \subset A_2/A_1 \subset A_3/A_1 \subset \ldots \subset A_n/A_1
$$
and
$$
0 \subset (B_1 + A_1)/A_1 \subset \ldots \subset
(B_{j - 1} + A_1)/A_1 = B_j/A_1 \subset B_{j + 1}/A_1
\subset \ldots \subset B_m/A_1
$$
We conclude by induction.
\end{proof}








\section{Serre subcategories}
\label{section-serre-subcategories}

\noindent
In \cite[Chapter I, Section 1]{Serre_homotopie_classes}
a notion of a ``class'' of abelian groups is defined.
This notion has been extended to abelian categories by many authors
(in slightly different ways). We will use the following variant
which is virtually identical to Serre's original definition.

\begin{definition}
\label{definition-serre-subcategory}
\begin{reference}
\cite[Condition (I) on page 259]{Serre_homotopie_classes}
\end{reference}
Let $\mathcal{A}$ be an abelian category.
\begin{enumerate}
\item A {\it Serre subcategory} of $\mathcal{A}$ is a
nonempty full subcategory $\mathcal{C}$ of $\mathcal{A}$
such that given an exact sequence\footnote{By
Definition \textnormal{[original label: \nolinkurl{definition-exact}]} this means $\Im(A \to B) = \Ker(B \to C)$.}
$$
A \to B \to C
$$
with $A, C \in \Ob(\mathcal{C})$, then also
$B \in \Ob(\mathcal{C})$.
\item A {\it weak Serre subcategory} of $\mathcal{A}$ is a nonempty
full subcategory $\mathcal{C}$ of $\mathcal{A}$ such that given an
exact sequence
$$
A_0 \to A_1 \to A_2 \to A_3 \to A_4
$$
with $A_0, A_1, A_3, A_4$ in $\mathcal{C}$, then also $A_2$ in $\mathcal{C}$.
\end{enumerate}
\end{definition}

\noindent
In some references the second notion is called a ``thick'' subcategory
and in other references the first notion is called a ``thick'' subcategory.
However, it seems that the notion of a Serre subcategory is universally
accepted to be the one defined above. Note that in both cases the category
$\mathcal{C}$ is abelian and that the inclusion functor
$\mathcal{C} \to \mathcal{A}$ is a fully faithful exact functor.
Let's characterize these types of subcategories in more detail.

\begin{lemma}
\label{lemma-characterize-serre-subcategory}
Let $\mathcal{A}$ be an abelian category.
Let $\mathcal{C}$ be a subcategory of $\mathcal{A}$.
Then $\mathcal{C}$ is a Serre subcategory if and only if
the following conditions are satisfied:
\begin{enumerate}
\item $0 \in \Ob(\mathcal{C})$,
\item $\mathcal{C}$ is a strictly full subcategory of $\mathcal{A}$,
\item any subobject or quotient of an object of $\mathcal{C}$ is an object
of $\mathcal{C}$,
\item if $A \in \Ob(\mathcal{A})$ is an extension of objects of $\mathcal{C}$
then also $A \in \Ob(\mathcal{C})$.
\end{enumerate}
Moreover, a Serre subcategory is an abelian category and
the inclusion functor is exact.
\end{lemma}

\begin{proof}
Omitted.
\end{proof}

\begin{lemma}
\label{lemma-characterize-weak-serre-subcategory}
Let $\mathcal{A}$ be an abelian category.
Let $\mathcal{C}$ be a subcategory of $\mathcal{A}$.
Then $\mathcal{C}$ is a weak Serre subcategory if and only if
the following conditions are satisfied:
\begin{enumerate}
\item $0 \in \Ob(\mathcal{C})$,
\item $\mathcal{C}$ is a strictly full subcategory of $\mathcal{A}$,
\item kernels and cokernels in $\mathcal{A}$ of morphisms
between objects of $\mathcal{C}$ are in $\mathcal{C}$,
\item if $A \in \Ob(\mathcal{A})$ is an extension of objects of $\mathcal{C}$
then also $A \in \Ob(\mathcal{C})$.
\end{enumerate}
Moreover, a weak Serre subcategory is an abelian category and
the inclusion functor is exact.
\end{lemma}

\begin{proof}
Omitted.
\end{proof}

\begin{lemma}
\label{lemma-kernel-exact-functor}
Let $\mathcal{A}$, $\mathcal{B}$ be abelian categories.
Let $F : \mathcal{A} \to \mathcal{B}$ be an exact functor.
Then the full subcategory of objects $C$ of $\mathcal{A}$
such that $F(C) = 0$ forms a Serre subcategory of $\mathcal{A}$.
\end{lemma}

\begin{proof}
Omitted.
\end{proof}

\begin{definition}
\label{definition-kernel-category}
Let $\mathcal{A}$, $\mathcal{B}$ be abelian categories.
Let $F : \mathcal{A} \to \mathcal{B}$ be an exact functor.
Then the full subcategory of objects $C$ of $\mathcal{A}$
such that $F(C) = 0$ is called the {\it kernel of the functor $F$},
and is sometimes denoted $\Ker(F)$.
\end{definition}

\noindent
Any Serre subcategory of an abelian category is the kernel of
an exact functor. In
Examples, Section \textnormal{[original label: \nolinkurl{examples-section-serre-quotient-modulo-torsion-modules}]}
we discuss this for Serre's original example of torsion groups.


\begin{lemma}
\label{lemma-serre-subcategory-is-kernel}
Let $\mathcal{A}$ be an abelian category.
Let $\mathcal{C} \subset \mathcal{A}$ be a Serre subcategory.
There exists an abelian category $\mathcal{A}/\mathcal{C}$
and an exact functor
$$
F : \mathcal{A} \longrightarrow \mathcal{A}/\mathcal{C}
$$
which is essentially surjective and whose kernel is $\mathcal{C}$.
The category $\mathcal{A}/\mathcal{C}$ and the functor $F$ are
characterized by the following universal property: For any exact
functor $G : \mathcal{A} \to \mathcal{B}$ such that
$\mathcal{C} \subset \Ker(G)$ there exists a factorization
$G = H \circ F$ for a unique exact functor
$H : \mathcal{A}/\mathcal{C} \to \mathcal{B}$.
\end{lemma}

\begin{proof}
Consider the set of arrows of $\mathcal{A}$ defined by
the following formula
$$
S = \{f \in \text{Arrows}(\mathcal{A}) \mid
\Ker(f), \Coker(f) \in \Ob(\mathcal{C}) \}.
$$
We claim that $S$ is a multiplicative system. To prove this we have
to check MS1, MS2, MS3, see
Categories, Definition \textnormal{[original label: \nolinkurl{categories-definition-multiplicative-system}]}.

\medskip\noindent
It is clear that identities are elements of $S$. Suppose that
$f : A \to B$ and $g : B \to C$ are elements of $S$.
There are exact sequences
$$
\begin{matrix}
0 \to \Ker(f) \to \Ker(gf) \to \Ker(g) \\
\Coker(f) \to \Coker(gf) \to \Coker(g) \to 0
\end{matrix}
$$
Hence it follows that $gf \in S$. This proves MS1. (In fact, a similar
argument will show that $S$ is a saturated multiplicative system, see
Categories, Definition
\textnormal{[original label: \nolinkurl{categories-definition-saturated-multiplicative-system}]}.)

\medskip\noindent
Consider a solid diagram
$$
\xymatrix{
A \ar[d]_t \ar[r]_g & B \ar@{..>}[d]^s \\
C \ar@{..>}[r]^f & C \amalg_A B
}
$$
with $t \in S$. Set
$W = C \amalg_A B =  \Coker((t, -g) : A \to C \oplus B)$.
Then $\Ker(t) \to \Ker(s)$ is surjective and
$\Coker(t) \to \Coker(s)$ is an isomorphism. Hence
$s$ is an element of $S$. This proves LMS2 and the proof of RMS2 is dual.

\medskip\noindent
Finally, consider morphisms $f, g : B \to C$ and a morphism $s : A \to B$
in $S$ such that $f \circ s = g \circ s$. This means that
$(f - g) \circ s = 0$. In turn this means that
$I = \Im(f - g) \subset C$ is a quotient of $\Coker(s)$
hence an object of $\mathcal{C}$. Thus $t : C \to C' = C/I$ is an
element of $S$ such that $t \circ (f - g) = 0$, i.e., such that
$t \circ f = t \circ g$. This proves LMS3 and the proof of
RMS3 is dual.

\medskip\noindent
Having proved that $S$ is a multiplicative system we set
$\mathcal{A}/\mathcal{C} = S^{-1}\mathcal{A}$, and we set
$F$ equal to the localization functor $Q$. By
Lemma \ref{lemma-localization-abelian}
the category $\mathcal{A}/\mathcal{C}$ is abelian and $F$ is exact.
If $X$ is in the kernel of $F = Q$, then by
Lemma \ref{lemma-kernel-localization}
we see that $0 : X \to Z$ is an element of $S$ and hence
$X$ is an object of $\mathcal{C}$, i.e., the kernel of
$F$ is $\mathcal{C}$.
Finally, if $G$ is as in the statement of the lemma, then $G$ turns
every element of $S$ into an isomorphism. Hence we obtain the
functor $H : \mathcal{A}/\mathcal{C} \to \mathcal{B}$ from
the universal property of localization, see
Categories, Lemma \textnormal{[original label: \nolinkurl{categories-lemma-properties-left-localization}]}.
We still have to show the functor $H$ is exact.
To do this it suffices to show that $H$ commutes
with taking kernels and cokernels, see Lemma \textnormal{[original label: \nolinkurl{lemma-exact-functor}]}.
Let $A \to B$ be a morphism in $\mathcal{A}/\mathcal{C}$.
We may represent $A \to B$ as $fs^{-1}$ where $s : A' \to A$
is in $S$ and $f : A' \to B$ an arbitrary morphism of $\mathcal{A}$.
Since $F = Q$ maps $s$ to an isomorphism in the quotient category
$\mathcal{A}/\mathcal{C}$, it suffices to show that $H$ commutes with taking
kernels and cokernels of morphisms $f : A \to B$ of $\mathcal{A}$.
But here we have $H(f) = G(f)$ and the result follows
from the fact that $G$ is exact.
\end{proof}


\begin{lemma}
\label{lemma-quotient-by-kernel-exact-functor}
Let $\mathcal{A}$, $\mathcal{B}$ be abelian categories.
Let $F : \mathcal{A} \to \mathcal{B}$ be an exact functor.
Let $\mathcal{C} \subset \mathcal{A}$ be a Serre subcategory
contained in the kernel of $F$.
Then $\mathcal{C} = \Ker(F)$ if and only if the induced functor
$\overline{F} : \mathcal{A}/\mathcal{C} \to \mathcal{B}$
(Lemma \ref{lemma-serre-subcategory-is-kernel}) is faithful.
\end{lemma}

\begin{proof}
We will use the results of Lemma \ref{lemma-serre-subcategory-is-kernel}
without further mention.
The ``only if'' direction is true because the kernel of $\overline{F}$ is zero
by construction. Namely, if $f : X \to Y$ is a morphism in
$\mathcal{A}/\mathcal{C}$ such that $\overline{F}(f) = 0$, then
$\overline{F}(\Im(f)) = \Im(\overline{F}(f)) = 0$, hence $\Im(f) = 0$ by the
assumption on the kernel of $F$. Thus $f = 0$.

\medskip\noindent
For the ``if'' direction, let $X$ be an object of $\mathcal{A}$ such that $F(X)
= 0$. Then $\overline{F}(\text{id}_X) = \text{id}_{\overline{F}(X)} = 0$, thus
$\text{id}_X = 0$ in $\mathcal{A}/\mathcal{C}$ by faithfulness of
$\overline{F}$. Hence $X = 0$ in $\mathcal{A}/\mathcal{C}$, that is $X \in
\Ob(\mathcal{C})$.
\end{proof}







\bibliographystyle{plain}
\bibliography{my}
\end{document}
